\thispagestyle{plain}
\phantomsection
\addcontentsline{toc}{chapter}{中文摘要}
\centerline{\zihao{-3}\heiti 中文摘要}

\linespread{1.4}\zihao{-4}\bigskip

针对地下停车场应急救援场景中 GNSS 信号不可用、环境多路径严重、人员轨迹连续性强等特点，本文研究“蓝牙信标 RSSI 观测 + 步行运动约束”的室内定位建模问题。为避免仅依赖 RSSI 反演导致的位置抖动，本文将信号距离模型与行人短时运动模型联合，构建分层优化框架：先由路径损耗模型给出粗定位，再利用带正则项的加权最小二乘进行轨迹平滑与实时修正。

在建模方面，本文建立了 RSSI 与距离的对数衰减关系，并引入异常抑制权重以减弱遮挡与反射造成的粗差影响；在约束方面，结合步长范围、转向速率上界与停车场车道拓扑约束，形成可解释的可行域。算法方面，采用“鲁棒初值 + 迭代重加权最小二乘 + 滑动窗口平滑”方案，兼顾精度、稳定性与计算效率。实验部分基于仿真停车场平面布局，设置多组信标密度、噪声水平和遮挡比例，对定位误差、轨迹连续性、鲁棒性及实时性进行系统评估。

结果表明：与仅使用 RSSI 的点位反演相比，融合模型在中高噪声场景下误差显著降低，轨迹抖动明显减小；当部分信标失效时，步行约束可有效维持解的稳定性；在 1 s 采样频率下，算法可满足应急定位的实时计算需求。本文形成了一套可直接迁移到地下停车场、地下商超及隧道巡检场景的建模流程，对应急人员定位具有实用参考价值。

\bigskip
\noindent{\zihao{4}\heiti 关键词：}
室内定位；地下停车场；RSSI；加权最小二乘；轨迹平滑；鲁棒估计
